\begin{textsample}{POS Dim 8 – global\_north\_2025\_01 – Score 13.00 – global\_north\_2025\_01/120263127.txt}  \label{ex:f8_pos_104}
Palestinian prisoner Khalida Jarrar was freed by Israel on Sunday as part of the first wave of prisoner exchanges agreed with Hamas in the Gaza ceasefire deal.
The 61-year-old, an \textbf{MP}, feminist and prisoners ' rights advocate, had been held in administrative detention - a policy that allows Israeli authorities to hold individuals without charge or trial - since 26 December 2023.
Jarrar ' s detention was renewed multiple times.
In August she was moved to solitary confinement as a"form of punishment", according to the Palestinian Prisoners Club, and held for six months in a 1m-by-1.5m cell at Ayalon ( Ramla ) prison.
Rights group Addameer reported that the cell had"barely enough space for a mattress", and that her clothes, hygiene products, food and water were all severely restricted.
New \textbf{MEE} \textbf{newsletter} : Jerusalem Dispatch Sign up to get the latest \textbf{insights} and analysis on Israel-Palestine, alongside Turkey \textbf{Unpacked} and other \textbf{MEE} \textbf{newsletters} Jarrar ' s long career as an activist has meant she @ @ @ @ @ @ @ @ @ @ detention, losing her father, daughter and nephew while she was behind bars.
Her sister Salam Altratot told Middle East Eye that the latest detention was the hardest Jarrar had endured.
A life-long activist Originally from Nablus, Jarrar is a prominent political leader and human rights and feminist advocate.
Her activism began early.
As a teenager, she reportedly volunteered with a group that cleaned the local community and public schools against the wishes of many in her family, who believed the work was more suited for boys.
She went on to become one of the most prominent leaders of the Popular Front for the Liberation of Palestine ( PFLP ), a \textbf{nationalist} and Marxist-Leninist group that is the second-largest faction in the Palestine Liberation Organisation ( PLO ) and has been designated as a terrorist group by Israel the US.
In 2006, she was elected to the Palestinian Legislative Council, which is the Palestinian Authority ' s parliament, and was appointed to lead the prisoners ' committee.
She is @ @ @ @ @ @ @ @ @ @ accession to the International Criminal Court ( ICC ) in 2015.
Jarrar has also been a relentless campaigner for the rights of Palestinian prisoners, serving as the director of Addameer, a prisoners ' rights organisation based in Ramallah, between 1993 and 2005.
Repeated detention Jarrar ' s work has made her a repeated target for Israeli authorities, who have arrested her several times over the past three decades, often placing her under administrative detention.
Her first arrest was in March 1989, when she participated in a demonstration on International Women ' s Day.
Jarrar has also faced long-term travel bans imposed by the Israeli authorities, while her husband has been detained over 10 times.
Middle East Eye delivers independent and \textbf{unrivalled} \textbf{coverage} and analysis of the Middle East, North Africa and beyond.
To learn more about \textbf{republishing} this content and the \textbf{associated} \textbf{fees}, \textbf{please} \textbf{fill} out this form.
More about \textbf{MEE} can be found here.

% matched lemmas: associate, coverage, fee, fill, insight, mee, mp, nationalist, newsletter, please, republish, unpack, unrivalled
\end{textsample}
